\subsection{Threats to validity}
\label{sec:threats}

\textbf{One machine, compressed time.} All runs used one 8-core desktop machine and a
surge compressed from tens of minutes to \SI{90}{s}. We therefore report load
as multiples of a calibrated capacity and compare designs relative to each
other; absolute rates do not transfer to production hardware. The
pre-computed designs were faster than our load generator could saturate
(C1 served at least \CapCOne\ students/s, the generator's limit), so their
margins are lower bounds.

\textbf{Modelled students.} No published study measures how students refresh
on result day. Patience, refresh probabilities and delays are assumptions;
Section~\ref{sec:refresh} shows how much the conclusions depend on them.

\textbf{Modelled server cost.} The \SI{5}{ms} page-building cost stands in for
the framework overhead of real portal stacks. With \SI{2}{ms} or \SI{0}{ms},
A and B0 took the 2$\times$ load easily (Section~\ref{sec:rq1}), so absolute
capacities depend on it; the cost order of the designs did not.

\textbf{Scope.} We do not model login or captcha services, bots, denial of
service, DNS, or a real CDN, each of which has caused real result-day outages.
Results are synthetic; three repetitions per configuration give wide
intervals for some metrics.
